
\documentclass[11pt,a4paper]{article}
\usepackage[utf8]{inputenc}
\usepackage{geometry}
\geometry{top=1in,bottom=1in,left=1in,right=1in}
\usepackage{amsmath,amssymb}
\usepackage{booktabs}
\usepackage{graphicx}
\usepackage{hyperref}
\usepackage{xcolor}
\usepackage{listings}
\usepackage{fancyhdr}

\pagestyle{fancy}
\fancyhf{}
\setlength{\headheight}{14pt}
\rhead{23CSE401 · Fundamentals of Artificial Intelligence · Assignment 2}
\lhead{Multi-Agent Air Traffic Conflict Resolution}
\cfoot{\thepage}

\title{\textbf{Multi-Agent Air Traffic Conflict Resolution and Route Planning}\\
\large Space-Time Heuristic Search, Prioritized Deconfliction, and Destination Parking Invariance in Congested Airspace\\
\vspace{0.3cm}
\normalsize Course: 23CSE401 — Fundamentals of Artificial Intelligence (Assignment 2)}
\author{\textbf{Academic Implementation \& Research Report}}
\date{October 16, 2026}

\begin{document}

\maketitle
\thispagestyle{empty}

\begin{abstract}
Air traffic management in terminal control areas requires autonomous multi-agent coordination capable of guaranteeing spatial and temporal separation under dynamic flight trajectories. In this project, we implement and empirically evaluate a complete, deterministic, space-time multi-agent pathfinding simulation. We formulate the navigation problem over an $(x, y, t)$ state-space representation, comparing \textbf{A* Search} (guided by an admissible Manhattan distance heuristic) against \textbf{Uniform-Cost Search (UCS)} as an unguided baseline. Multi-agent interactions are coordinated through prioritized space-time reservation tables equipped with transition locks that explicitly eliminate vertex collisions, head-on edge-swaps, and destination parking conflicts. When bottlenecks emerge, agents dynamically execute in-air holding maneuvers (wait actions) or bounded priority-reordering replanning. Experimental results across seven benchmark scenarios reveal that A* matches UCS's per-aircraft path cost under identical reservations while expanding up to 96\% fewer nodes. Because planning is prioritized, total cost across all aircraft is not guaranteed to be globally minimal (Scenario 7: A* 90.0 vs UCS 89.0). Both algorithms demonstrate robust safety, successfully proving infeasibility by exhausting the reachable space-time graph within the planning horizon and maintaining zero unresolved conflicts across all executed benchmark trajectories verified by an independent safety monitor.
\end{abstract}

\vspace{0.5cm}
\tableofcontents
\newpage

\section{Introduction}
Modern air traffic control (ATC) is transitioning from human-centric sector management to autonomous NextGen operations capable of handling high-density commercial airliners, unmanned aerial vehicles (UAVs), and urban air mobility (UAM) shuttles. The core challenge of autonomous ATC is multi-agent collision avoidance: computing flight trajectories that safely separate aircraft while minimizing transit delay and fuel consumption.

Independent single-agent pathfinding algorithms frequently produce fatal schedule conflicts because shortest spatial paths often intersect temporally. Conversely, fully coupled multi-agent state space representations grow exponentially ($O(b^k)$) with the number of aircraft $k$, rendering exhaustive search intractable. In this paper, we bridge this trade-off using prioritized space-time planning \cite{silver2005}, implementing genuine priority-queue search engines, reservation matrices, and an interactive simulation dashboard.

\section{Problem Domain and Formal Statement}
Following standard multi-agent pathfinding (MAPF) formalisms \cite{stern2019}, consider an airspace discretized into a two-dimensional grid of dimensions $W \times H$, populated by static no-fly zones (restricted cells $R \subset \mathbb{Z}^2$) and a set of $k$ aircraft agents $\mathcal{A} = \{A_1, A_2, \dots, A_k\}$. Each aircraft $A_i$ is characterized by an initial position $s_i \in \mathbb{Z}^2 \setminus R$, a destination goal $g_i \in \mathbb{Z}^2 \setminus R$, an assigned 0-based priority rank $p_i \in \{0, 1, \dots, k-1\}$ (where lower numerical value denotes higher priority, e.g., $p=0$ for $A_1$), and unit transit costs.

The objective is to compute a collection of discrete-time collision-free trajectories $\Pi = \{\pi_1, \pi_2, \dots, \pi_k\}$, where trajectory $\pi_i = \langle (x_i^0, y_i^0), (x_i^1, y_i^1), \dots, (x_i^{T_i}, y_i^{T_i}) \rangle$ satisfies:
\begin{enumerate}
    \item \textbf{Endpoint Integrity:} $(x_i^0, y_i^0) = s_i$ and $(x_i^{T_i}, y_i^{T_i}) = g_i$.
    \item \textbf{Kinematic Validity:} For all time steps $t$, $\|(x_i^{t+1}, y_i^{t+1}) - (x_i^t, y_i^t)\|_1 \le 1$. The discrete action set is $\{\text{Up, Down, Left, Right, Wait}\}$.
    \item \textbf{Static Obstacle Clearance:} For all $t$, $(x_i^t, y_i^t) \notin R$.
    \item \textbf{Vertex Separation:} $\forall i \neq j, \forall t, (x_i^t, y_i^t) \neq (x_j^t, y_j^t)$.
    \item \textbf{Edge-Swap Separation:} $\forall i \neq j, \forall t$, no head-on swap: $((x_i^t, y_i^t) = (x_j^{t+1}, y_j^{t+1})) \land ((x_i^{t+1}, y_i^{t+1}) = (x_j^t, y_j^t))$.
    \item \textbf{Destination Parking Invariance:} When aircraft $A_i$ reaches goal $g_i$ at time $T_i$, it remains parked at $g_i$ for all subsequent time steps $t \in [T_i, T_{\max}]$. Thus, $\forall j \neq i, \forall t \ge T_i, (x_j^t, y_j^t) \neq g_i$.
    \item \textbf{Cost Minimization:} Minimize individual route cost:
    \[
    C(\pi_i) = \sum_{t=0}^{T_i-1} c\big((x_i^t, y_i^t), (x_i^{t+1}, y_i^{t+1})\big), \quad \text{where unit transit/wait cost } c(\cdot) = 1.0.
    \]
\end{enumerate}

\subsection{Benchmark Scenario Definitions}
Table~\ref{tab:scenarios} details the seven standard benchmark scenarios utilized throughout this investigation.

\begin{table}[htbp]
\centering
\caption{Deterministic Benchmark Scenario Specifications}
\label{tab:scenarios}
\resizebox{\textwidth}{!}{%
\begin{tabular}{lccll}
\toprule
\textbf{Scenario} & \textbf{Grid ($W \times H$)} & \textbf{Agents} & \textbf{Aircraft Configurations: ID (Start $\to$ Goal, Priority)} & \textbf{Obstacles / Restrictions} \\
\midrule
1. Conflict-Free & $10 \times 10$ & 2 & AC-1: $(0, 2) \to (9, 2), p=0$; AC-2: $(0, 7) \to (9, 7), p=1$ & None \\
2. Crossing Routes & $8 \times 8$ & 2 & AC-1: $(0, 3) \to (7, 3), p=0$; AC-2: $(3, 0) \to (3, 7), p=1$ & None (Crosses at $(3, 3)$) \\
3. Unreachable Destination & $10 \times 10$ & 1 & AC-1: $(0, 0) \to (5, 4), p=0$ & 8 cells enclosing $(5, 4)$ \\
4. Narrow Corridor & $15 \times 7$ & 3 & AC-1: $(0, 3) \to (14, 3), p=0$; AC-2: $(1, 3) \to (13, 3), p=1$; AC-3: $(2, 3) \to (12, 3), p=2$ & Walls at $y \in \{2, 4\}$ forming corridor \\
5. Start Equals Goal & $5 \times 5$ & 1 & AC-1: $(2, 2) \to (2, 2), p=0$ & None \\
6. Head-On Collision & $10 \times 10$ & 2 & AC-1: $(1, 5) \to (8, 5), p=0$; AC-2: $(8, 5) \to (1, 5), p=1$ & None \\
7. Five Aircraft (Default) & $12 \times 12$ & 5 & AC-1: $(0, 2) \to (11, 2), p=0$; AC-2: $(11, 4) \to (0, 4), p=1$; AC-3: $(5, 0) \to (5, 11), p=2$; & 3 static blockades \\
 & & & AC-4: $(2, 10) \to (9, 3), p=3$; AC-5: $(9, 10) \to (2, 3), p=4$ & total 10 cells \\
\bottomrule
\end{tabular}%
}
\end{table}

\section{PEAS Framework}
The air traffic coordination system is formalized under Russell and Norvig's PEAS framework \cite{russell2020}:
\begin{itemize}
    \item \textbf{Performance Measure:} Zero vertex and edge collisions, zero restricted zone penetrations, minimal cumulative path cost $\sum C(\pi_i)$, minimal waiting delay $\sum W_i$, minimized makespan $\max_i(T_i)$, and mission completion for all feasible, non-deadlocking flight configurations (subject to the known incompleteness of decoupled prioritized planning).
    \item \textbf{Environment:} Discrete 2D grid airspace with static restricted zones, discrete unit time clock $t \in [0, T_{\max}]$, dynamic reservation tables, and deterministic transition physics.
    \item \textbf{Actuators:} Grid navigation commands: $\Delta x \in \{-1, 0, 1\}$, $\Delta y \in \{-1, 0, 1\}$ subject to $\|\Delta x\| + \|\Delta y\| \le 1$.
    \item \textbf{Sensors:} Global coordinates $(x, y)$, current time step $t$, goal coordinate $(g_x, g_y)$, and active space-time reservation queries.
\end{itemize}

\section{Environment Properties}
Following \cite{russell2020}, the simulation airspace exhibits the following core AI environment characteristics:
\begin{enumerate}
    \item \textbf{Fully Observable:} The centralized coordinator possesses complete, omniscient knowledge of the grid geometry, obstacle placements, aircraft coordinates, and schedule reservations across all time horizons.
    \item \textbf{Deterministic:} Actions yield strictly predictable outcomes without stochastic slip, sensor noise, or wind drift.
    \item \textbf{Sequential:} A decision made at time $t$ commits an airspace reservation that constrains all future trajectories for lower-priority agents.
    \item \textbf{Static:} Planning is conducted offline over a static physical environment with static no-fly zones and deterministic transitions, while time-indexed reservations evolve deterministically.
    \item \textbf{Discrete:} Coordinates, time steps, and action choices are integer-valued.
    \item \textbf{Multi-Agent Cooperative:} Agents collaborate via centralized prioritized arbitration to achieve system-wide safety and throughput.
\end{enumerate}

\section{Agent Architecture and Cooperative Strategy}
We implement a \textbf{Goal-Based Cooperative Agent Architecture}. Aircraft agents are modeled as goal-based entities: their behavior is guided strictly by reaching destination coordinates $g_i$ while minimizing transit cost. They do not employ a multi-attribute utility function; instead, coordination follows \textbf{Prioritized Space-Time Planning} \cite{silver2005}:
\begin{enumerate}
    \item Aircraft are ranked deterministically by assigned priority rank $p_i$, breaking ties deterministically by unique aircraft ID.
    \item The highest-priority aircraft $A_1$ plans an optimal path to its goal ignoring other agents, but respecting static obstacles.
    \item The planned trajectory of $A_1$ is committed to the space-time reservation table $\mathcal{R}$, including destination parking through $T_{\max}$.
    \item Successive aircraft $A_2, \dots, A_k$ plan space-time paths treating prior reservations in $\mathcal{R}$ as transient obstacles.
    \item \textbf{Destination Parking Invariant:} When planning a path for aircraft $A_i$ arriving at goal $g_i$ at time $t$, the search algorithm verifies that $g_i$ is free in $\mathcal{R}$ for all $t' \in [t, T_{\max}]$. Once confirmed, $\mathcal{R}$ reserves $(g_i, t')$ for all $t' \ge t$, ensuring lower-priority flights cannot collide with a parked aircraft.
    \item If an agent cannot reach its goal within the horizon $T_{\max}$, or if an independent conflict monitor detects residual separation violations, the coordinator initiates \textbf{Bounded Replanning} by rotating priority rankings.
    \item \textbf{Partial Planning and Fault Containment:} If an individual aircraft cannot reach its destination (e.g., due to an enclosed goal), but the remaining aircraft have planned paths with zero mutual conflicts, the coordinator allows the conflict-free flights to execute while safely holding the failed aircraft parked at its departure point for all $t \in [0, T_{\max}]$. If any unresolved conflicts remain among flights, the simulation strictly halts without execution.
\end{enumerate}

\section{Search Algorithms: A* vs Uniform-Cost Search}
\subsection{A* Search}
A* \cite{hart1968} expands nodes in order of the evaluation function:
\begin{equation}
f(n) = g(n) + h(n)
\end{equation}
where $g(n)$ is the exact accumulated path cost from the start configuration $(x_0, y_0, 0)$ to state $n = (x, y, t)$, and $h(n)$ is the heuristic estimate to the destination $(g_x, g_y)$.

We employ the \textbf{Manhattan Distance Heuristic}:
\begin{equation}
h(x, y) = |x - g_x| + |y - g_y|
\end{equation}
\textbf{Admissibility Proof:} Since all cardinal grid transitions incur a unit cost of $c=1.0$, diagonal transitions are forbidden, and in-air waiting incurs a cost of $1.0$, the minimum possible cost between any two points in the absence of obstacles is exactly the Manhattan distance. Therefore, $h(n) \le h^*(n)$ for all states $n$. With an admissible (and here consistent) heuristic, A* returns a minimum-cost path for a single aircraft given a fixed reservation table, non-negative costs and the horizon $T_{\max}$. This is a per-search guarantee; it does not extend to the joint cost of all aircraft.

\subsection{Uniform-Cost Search (UCS)}
Uniform-Cost Search acts as the baseline unguided search algorithm. It maintains a priority queue ordered strictly by $g(n)$, equivalent to setting $h(n) = 0$. While UCS guarantees cost optimality under non-negative transition costs for an individual search under fixed reservations, it explores states isotropically in all directions.

\section{State-Space Representation and Search Tree}
A state in the space-time search graph is represented by the 3-tuple:
\begin{equation}
s = (x, y, t) \in \{0, \dots, W-1\} \times \{0, \dots, H-1\} \times \{0, \dots, T_{\max}\}
\end{equation}
For a current state $s = (x, y, t)$, the valid successor states $s' = (x', y', t+1)$ satisfy:
\begin{equation}
(x', y') \in \{(x+1, y), (x-1, y), (x, y+1), (x, y-1), (x, y)\} \setminus R
\end{equation}
subject to reservation table constraints:
\begin{equation}
(x', y', t+1) \notin \mathcal{R}_{\text{cells}} \quad \text{and} \quad ((x, y, t), (x', y', t+1)) \notin \mathcal{R}_{\text{edges}}
\end{equation}
The branching factor $b = 5$ (four cardinal moves plus wait). The size of the space-time search graph is bounded by $O(W \cdot H \cdot T_{\max})$ distinct states.

\section{Sequence of Explored States: Real Measured Trace}
To verify the internal mechanics of the search algorithm, the engine records an exact expansion trace (\texttt{ExpansionRecord}) whenever tracing is enabled. Each record logs the expansion step, space-time state $(x, y, t)$, $g(n)$, $h(n)$, $f(n)$, and the count of successors blocked by the reservation table.

\subsection{Scenario 2: Crossing Routes Exploration Trace}
In Scenario 2, Aircraft AC-1 departs $(0, 3)$ bound for $(7, 3)$, while AC-2 departs $(3, 0)$ bound for $(3, 7)$. AC-1 plans first and reserves the horizontal row $y=3$, occupying intersection cell $(3, 3)$ at $t=3$.

When AC-2 (Priority 1) plans its route from $(3, 0)$ to $(3, 7)$ under A*, the priority queue expands the sequence of states shown in Table~\ref{tab:trace_ac2}.

% Auto-generated by generate_trace.py — DO NOT EDIT MANUALLY
\begin{table}[htbp]
\centering
\caption{Real Measured Expansion Trace for AC-2 in Crossing Routes Scenario (A*). Notes column is descriptive text added for Scenario 2.}
\label{tab:trace_ac2}
\resizebox{\textwidth}{!}{%
\begin{tabular}{ccccccc}
\toprule
\textbf{Step} & \textbf{State $(x, y, t)$} & \textbf{$g(n)$} & \textbf{$h(n)$} & \textbf{$f(n)$} & \textbf{Blocked} & \textbf{Search Action / Notes} \\
\midrule
1 & $(3, 0, t=0)$ & 0.0 & 7.0 & 7.0 & 0 & Start state popped from open list \\
2 & $(3, 1, t=1)$ & 1.0 & 6.0 & 7.0 & 0 & Unconstrained move towards destination \\
3 & $(3, 2, t=2)$ & 2.0 & 5.0 & 7.0 & 1 & Adjacent to crossing; successor $(3, 3, 3)$ blocked by AC-1 reservation (1 pruned) \\
4 & $(3, 2, t=3)$ & 3.0 & 5.0 & 8.0 & 0 & In-air holding maneuver (wait state at $h=5.0$, $f=8.0$) \\
5 & $(3, 3, t=4)$ & 4.0 & 4.0 & 8.0 & 0 & Advances to crossing cell $(3, 3)$ after AC-1 clears \\
6 & $(3, 4, t=5)$ & 5.0 & 3.0 & 8.0 & 0 & Unconstrained move towards destination \\
7 & $(3, 5, t=6)$ & 6.0 & 2.0 & 8.0 & 0 & Unconstrained move towards destination \\
8 & $(3, 6, t=7)$ & 7.0 & 1.0 & 8.0 & 0 & Unconstrained move towards destination \\
\midrule
\multicolumn{7}{c}{\textit{Goal $(3, 7, t=8)$ reached at $f=8.0$; popped and validated without further expansion (Total expansions: 8)}} \\
\bottomrule
\end{tabular}%
}
\end{table}


As recorded in the real measured trace, AC-2 advances unhindered from $(3, 0, 0)$ through $(3, 2, 2)$ (Steps 1--3) with $f=7.0$. At Step 3, the direct successor $(3, 3, 3)$ is blocked by AC-1's reservation and is pruned by the reservation table. The search then expands the in-air wait state at $(3, 2, 3)$ at Step 4 with $g=3.0, h=5.0, f=8.0$. After AC-1 clears the crossing, Step 5 advances into $(3, 3, 4)$ with $f=8.0$. Steps 6--8 climb vertically to $(3, 6, 7)$. The goal state $(3, 7, 8)$ is generated at $f=8.0$ and popped at the subsequent iteration where the destination parking invariant validates it, terminating the search with exactly 8 expansions and zero conflicts.

\section{Theoretical Complexity Analysis}
\begin{table}[ht]
\centering
\caption{Theoretical Complexity Breakdown}
\label{tab:complexity}
\begin{tabular}{lll}
\toprule
\textbf{Component} & \textbf{Time Complexity} & \textbf{Space Complexity} \\
\midrule
Single A* Search & $O(b^d)$ worst-case & $O(b^d)$ nodes in memory \\
Single UCS Search & $O(b^{C^*/\epsilon})$ & $O(b^{C^*/\epsilon})$ \\
Space-Time Grid Bound & $O(W \cdot H \cdot T_{\max})$ & $O(W \cdot H \cdot T_{\max})$ \\
Pairwise Conflict Detection & $O(k^2 \cdot T)$ & $O(k \cdot T)$ trajectory buffer \\
Prioritized Coordination ($k$ agents) & $O(k \cdot \text{Search}(W, H, T))$ & $O(k \cdot T)$ reservations \\
Bounded Replanning ($R$ attempts) & $O(R \cdot k \cdot \text{Search}(W, H, T))$ & $O(k \cdot T)$ \\
\bottomrule
\end{tabular}
\end{table}
The branching factor $b = 5$ (four cardinal moves plus wait). While worst-case tree search time is exponential in horizon depth $d$, the state-space bound for graph search is strictly $O(W \cdot H \cdot T_{\max})$ states. The admissible Manhattan distance heuristic drastically reduces the effective branching factor $b^*$ to near $1.0$ along unconstrained corridors.

\section{Empirical Experimental Results}
Experiments were executed on an x86\_64 architecture running Python 3.12 under Windows 11. Node counts, path costs, and expansion traces are strictly deterministic across runs, while wall-clock runtimes vary slightly with CPU operating state. Table~\ref{tab:comparison} presents the exact empirical results across all seven deterministic benchmark scenarios.

\begin{table}[htbp]
\centering
\caption{Empirical Comparison between A* and UCS across Benchmark Scenarios}
\label{tab:comparison}
\resizebox{\textwidth}{!}{%
\begin{tabular}{lcrrrrrrrrr}
\toprule
\textbf{Scenario} & \textbf{Algo} & \textbf{Success} & \textbf{Cost} & \textbf{Nodes Exp.} & \textbf{Nodes Gen.} & \textbf{Time (ms)} & \textbf{Wait} & \textbf{Makespan} & \textbf{Conflicts Before} & \textbf{Unresolved} \\
\midrule
1. Conflict-Free Routes & A* & True & 18.0 & 18 & 90 & 0.288 & 0 & 9 & 0 & \textbf{0} \\
1. Conflict-Free Routes & UCS & True & 18.0 & 456 & 594 & 4.280 & 0 & 9 & 0 & \textbf{0} \\
\midrule
2. Crossing Routes & A* & True & 15.0 & 15 & 74 & 0.433 & 1 & 8 & 1 & \textbf{0} \\
2. Crossing Routes & UCS & True & 15.0 & 334 & 439 & 3.964 & 1 & 8 & 1 & \textbf{0} \\
\midrule
3. Unreachable Destination & A* & False & N/A & 8,372 & 8,372 & 92.015 & N/A & N/A & 0 & \textbf{0} \\
3. Unreachable Destination & UCS & False & N/A & 8,372 & 8,372 & 83.400 & N/A & N/A & 0 & \textbf{0} \\
\midrule
4. Narrow Corridor & A* & True & 47.0 & 317 & 468 & 3.417 & 16 & 22 & 3 & \textbf{0} \\
4. Narrow Corridor & UCS & True & 47.0 & 974 & 1,097 & 8.697 & 2 & 22 & 3 & \textbf{0} \\
\midrule
5. Start Equals Goal & A* & True & 0.0 & 0 & 1 & 0.024 & 0 & 0 & 0 & \textbf{0} \\
5. Start Equals Goal & UCS & True & 0.0 & 0 & 1 & 0.032 & 0 & 0 & 0 & \textbf{0} \\
\midrule
6. Head-On Collision & A* & True & 16.0 & 20 & 87 & 0.282 & 0 & 9 & 1 & \textbf{0} \\
6. Head-On Collision & UCS & True & 16.0 & 388 & 504 & 4.665 & 0 & 9 & 1 & \textbf{0} \\
\midrule
7. Five Aircraft (Default) & A* & True & 90.0 & 336 & 761 & 5.343 & 4 & 22 & 3 & \textbf{0} \\
7. Five Aircraft (Default) & UCS & True & 89.0 & 6,117 & 6,807 & 69.743 & 1 & 22 & 3 & \textbf{0} \\
\bottomrule
\end{tabular}%
}
\vspace{0.2cm}
\footnotesize{\textit{Note: N/A indicates no feasible route exists for any aircraft (aircraft failed); 0.0 in Scenario 5 represents a genuine zero-cost route where start equals destination.}}
\end{table}

\section{Discussion of Findings}
\begin{enumerate}
    \item \textbf{Heuristic Pruning Efficiency:} In Scenario 1, A* expanded only 18 nodes compared to UCS's 456 nodes—a 96\% reduction in expanded states. In the dense 5-aircraft benchmark, A* expanded 336 nodes vs 6,117 nodes for UCS (18.2$\times$ fewer node expansions, translating to an empirical runtime speedup of approximately $13\times$), demonstrating practical viability in congested terminal airspaces.
    \item \textbf{Cost Optimality and Decoupled Planning Nuance:} In Scenarios 1 to 6, A* and UCS produced mathematically identical path costs (18.0, 15.0, N/A, 47.0, 0.0, and 16.0). Equal costs per search are consistent with heuristic admissibility, formally proven by $h(n) \le h^*(n)$ for Manhattan distance under unit edge costs. In Scenario 7 (5 Aircraft), A* achieved a total cost of 90.0 while UCS achieved 89.0. This discrepancy is a classical property of decoupled prioritized planning: when higher-priority agents encounter multiple co-optimal paths, A* and UCS break ties differently. These distinct tie-breaking choices commit different space-time reservations into table $\mathcal{R}$, which subsequently forces lower-priority agents into different detours. A* guarantees single-agent optimality given fixed reservations, but does not dominate UCS on total joint cost across all aircraft under prioritized planning.
    \item \textbf{Corridor Bottleneck Queuing and Wait Distributions:} In Scenario 4 (Narrow Corridor), both algorithms achieved the identical minimal route cost of 47.0 and identical makespan of 22 steps. However, A* incurred 16 in-air wait steps against UCS's 2 wait steps. This difference is a tie-breaking artifact: because the cost model charges in-air waiting and forward motion equally ($c=1.0$), trajectories with equal total length have identical cost. In a real aviation domain, holding patterns incur distinct operational and fuel penalties compared to linear detours.
    \item \textbf{Destination Parking Invariant and Safety Guarantee:} Reserving goal cells through $T_{\max}$ and verifying destination availability during search completely eliminated collision risks at destinations, guaranteeing \textbf{0 unresolved conflicts across all benchmark scenarios}.
    \item \textbf{Prioritized Planning vs Conflict-Based Search (CBS):} While complete algorithms like Conflict-Based Search (CBS) \cite{sharon2015} guarantee joint global optimality, CBS requires a two-level search with an exponential constraint tree ($O(2^C)$) that can timeout in dense terminal airspaces. Prioritized planning with bounded replanning operates in $O(R \cdot k \cdot \text{Search})$ polynomial time, providing predictable response latencies essential for tactical air traffic management.
    \item \textbf{Safe Handling of Infeasibility:} In Scenario 3, restricted cells completely enclosed the destination. Neither algorithm fabricated an illegal path; both safely exhausted the reachable space-time graph within the planning horizon (8,372 states) and returned a structured failure code with cost recorded as N/A.
    \item \textbf{Conflict Metrics Interpretation:} The ``Conflicts Before Coordination'' metric counts pairwise vertex and edge-swap collisions per time step among independent, uncoordinated single-agent baseline trajectories. The ``Unresolved Conflicts After Coordination'' metric counts separation violations remaining after planning. Across all benchmarks, the coordinator successfully reduces all pre-coordination conflicts to exactly 0 unresolved conflicts in the final coordinated plan.
\end{enumerate}

\section{System Architecture and Implementation}
The codebase is structured into modular Python packages adhering to rigorous separation of concerns:
\begin{itemize}
    \item \texttt{src/models.py}: Core data structures (\texttt{Aircraft}, \texttt{GridConfig}, \texttt{SearchResult}, \texttt{ExpansionRecord}, \texttt{Conflict}, etc.).
    \item \texttt{src/environment.py}: Grid boundary validation, neighborhood expansion ($b=5$), and movement cost calculus.
    \item \texttt{src/search.py}: Genuine priority-queue implementations of A* and UCS in space-time $(x, y, t)$ with destination parking checks and expansion tracing.
    \item \texttt{src/reservations.py}: Space-time reservation table with cell, edge, reverse-edge, and destination parking locks.
    \item \texttt{src/conflict\_detection.py}: Independent trajectory validator detecting vertex, edge-swap, restricted cells, and invalid movements.
    \item \texttt{src/coordinator.py}: Prioritized multi-agent coordinator with deterministic tie-breaking and bounded replanning.
    \item \texttt{src/simulation.py}: Lifecycle state machine and discrete-time simulation stepper.
    \item \texttt{src/metrics.py}: Metrics calculation, DataFrame conversion, and multi-scenario comparison table generation.
    \item \texttt{src/scenarios.py}: Programmatic scenario library and JSON loader.
    \item \texttt{generate\_outputs.py}: Automated batch runner generating benchmark CSV outputs.
    \item \texttt{generate\_trace.py}: Script generating the real expansion trace CSV and LaTeX table.
    \item \texttt{data/scenarios.json}: Machine-readable scenario definitions.
    \item \texttt{app.py}: Full Streamlit reactive dashboard with live grid animations, metrics display, and CSV export.
    \item \texttt{tests/test\_core.py}: 67 automated unit and integration tests passing with 100\% success rate.
\end{itemize}

\section{Limitations}
\begin{enumerate}
    \item \textbf{Suboptimality and Incompleteness of Prioritized Planning:} Prioritized planning is fundamentally incomplete and cannot guarantee joint cost optimality across all agents. In complex interlocking geometries, higher-priority aircraft may monopolize bottlenecks, causing lower-priority agents to fail even when joint conflict-free paths exist. Conflict-Based Search (CBS) \cite{sharon2015} provides a complete and optimal alternative.
    \item \textbf{Dependence on Priority Order:} Solution quality and feasibility depend strongly on the assigned priority order $p_i$, which bounded replanning can only partially mitigate via rotation heuristics.
    \item \textbf{Stationary Parking of Failed Aircraft:} When an aircraft's route planning fails, the system parks it at its departure cell for the entire horizon $T_{\max}$, permanently blocking that cell.
    \item \textbf{Centralized Architecture:} The system relies on a centralized coordinator representing a single point of failure.
    \item \textbf{Uniform Cost Model:} Waiting in place is charged identically to forward flight ($c=1.0$), abstracting the fuel and operational penalties of actual airborne holding patterns.
    \item \textbf{Discrete Kinematic Abstraction and Scale:} Continuous aircraft dynamics (turn radii, speed variation, wind vectors) are abstracted into unit grid steps, and scalability remains untested beyond 5 aircraft on larger grids.
\end{enumerate}

\section{Conclusion and Future Work}
This project successfully demonstrated the application of classical artificial intelligence search and multi-agent coordination to air traffic conflict resolution. A* proved dramatically superior to Uniform-Cost Search in computational efficiency (expanding up to 96\% fewer nodes) while preserving per-aircraft optimality under fixed reservations. The prioritized space-time reservation model with destination parking invariance eliminated all vertex and edge-swap collisions across seven benchmark scenarios and the automated test suite.

Future work includes exploring complete solvers such as Conflict-Based Search (CBS) \cite{sharon2015} or M*, incorporating continuous aircraft kinematics (Dubins curves), and deploying decentralized negotiation protocols.

\begin{thebibliography}{9}
\bibitem{hart1968}
P.~E. Hart, N.~J. Nilsson, and B.~Raphael.
\newblock A formal basis for the heuristic determination of minimum cost paths.
\newblock {\em IEEE Transactions on Systems Science and Cybernetics}, 4(2):100--107, 1968.

\bibitem{russell2020}
S.~Russell and P.~Norvig.
\newblock {\em Artificial Intelligence: A Modern Approach}.
\newblock Pearson, 4th edition, 2020.

\bibitem{silver2005}
D.~Silver.
\newblock Cooperative pathfinding.
\newblock In {\em Proceedings of the AAAI Conference on Artificial Intelligence and Interactive Digital Entertainment (AIIDE)}, volume~1, pages 117--122, 2005.

\bibitem{stern2019}
R.~Stern, N.~Sturtevant, A.~Felner, S.~Koenig, H.~Ma, T.~K. Walker, J.~Li, D.~Atzmon, L.~Cohen, T.~K. Kumar, et~al.
\newblock Multi-agent pathfinding: Definitions, variants, and benchmarks.
\newblock In {\em Twelfth Annual Symposium on Combinatorial Search (SoCS)}, 2019.

\bibitem{sharon2015}
G.~Sharon, R.~Stern, A.~Felner, and N.~R. Sturtevant.
\newblock Conflict-based search for optimal multi-agent pathfinding.
\newblock {\em Artificial Intelligence}, 219:40--66, 2015.
\end{thebibliography}

\end{document}